% pivot_entry_lemma.tex -- the entry lemma of the pivot band
% (round 382, LEMMA.PIVOT_ENTRY.01).  Builds with pdflatex.
% Companion machine check: verify_pivot_entry.py.
% Census probe: experiments/tfpt-discovery/pivot_entry_lemma_probe.py.
% NO Lean edit this round (R384 formalizes the named Prop).
% NO RH CLAIM.  Research documentation, not a theorem of RH.
\documentclass[11pt]{article}

\usepackage[a4paper,margin=2.7cm]{geometry}
\usepackage{amsmath,amssymb,amsthm}
\usepackage{microtype}
\usepackage{booktabs}
\usepackage{xcolor}
\usepackage[colorlinks=true,linkcolor=blue!50!black,urlcolor=blue!50!black]{hyperref}

\newtheorem{lemma}{Lemma}
\newtheorem{prop}{Proposition}
\newtheorem{definition}{Definition}
\theoremstyle{remark}
\newtheorem{remark}{Remark}

\newcommand{\R}{\mathbb{R}}
\newcommand{\indm}{\operatorname{ind}_{-}}
\newcommand{\In}{\operatorname{In}}

\title{\bfseries The entry lemma of the pivot band\\[2mm]
\large Flank interlacing, a source-defined $n_0=\lfloor 2N/5\rfloor$,
and the $L^2$ remainder that is half-filling}
\author{Stefan Hamann}
\date{August 28, 2026}

\begin{document}
\maketitle

\begin{abstract}\noindent
Round 380 proved that a rank-one Loewner chain cannot raise negative
inertia: an entry $\indm(A_{n_0})\le 1$ at a source-defined index
$n_0$ is transported through the whole band
(\texttt{adaptive\_band\_from\_entry}).  Equivalently, by the Jacobi
count, the monic pivots $h_0,\ldots,h_{n_0}$ of the signed source
Hankel have at most one sign change.  The strongest form is
$n_0=N-1$ with zero changes --- the free-window floor plate, which
is lemma $L^*$ itself.  The three-term recurrence alone does not
forbid a second prefix return (a matched scramble dies at degree
$21$, with $37$ negative pivots in the prefix).  This note isolates
the source ingredient that does: every negative run on the ordered
cosine support has length at most $2$, and the run's $\nu$-mass is
at most two-thirds of the immediately flanking $\mu$-mass.  Under
that condition the two-period adversary of amplitude $2/3$ first
goes negative at $\sim N/2$, and the reduced entry
$n_0=\lfloor 2N/5\rfloor$ sits strictly below that kill and strictly
above the scramble death.  What remains for $n_0=N-1$ is an $L^2$
restriction estimate for the Christoffel--Darboux kernel on the
flank blocks --- the same coherence remainder that $L^*$ is.
Nothing in this document is evidence for or against the Riemann
hypothesis.
\end{abstract}

\medskip\noindent
\textbf{Status.}\enspace
Lemmas~\ref{lem:rkhs}--\ref{lem:christoffel} are proved (finite
RKHS and Hankel algebra).
Lemmas~\ref{lem:toy-int}--\ref{lem:toy-cluster} are proved
(exact five-, three-, six- and seven-atom identities).
Lemma~\ref{lem:adv-tp} is proved (two-period $c\ge 1$ kills
$h_0$ or $H_1$).
Proposition~\ref{prop:census} is a finite machine census of $85$
MAIN windows, six terminal-dead $\chi$-sprouts, and a matched
scramble, not a cofinal law.
Proposition~\ref{prop:twoperiod} is a finite table of two-period
first negatives, not a theorem of every amplitude.
The original statement --- a source-defined $n_0\le N-1$ with
$\indm H_{n_0+1}\le 1$ on every canonical window --- is
\emph{reduced} to the flank condition of~\S\ref{sec:reduction}
plus the named $L^2$ remainder of~\S\ref{sec:half}.
No RH claim.  No $L^*$ claim.

\section{The object}\label{sec:obj}

Let $X=\{x_j=\cos(\pi j/S):j=1,\ldots,S\}$ with $S=2N-1$ odd, and
let $w$ be a real atomic weight on $X$ (the folded signed
von-Mangoldt tent density times the mesh factor $1-x$).  Write
$\mu$ for the positive part and $\nu$ for the negative part, so
$w=\mu-\nu$ as signed measures on the same grid, with disjoint
supports.  Write $H_r(\rho)$ for the $r\times r$ Hankel determinant
of a weight $\rho$, $H_0:=1$, and
\[
h_j(\rho)=\frac{H_{j+1}(\rho)}{H_j(\rho)}
\]
for the monic Jacobi norms whenever the minors are nonzero.  The
$\mu$-orthonormal Christoffel--Darboux kernel of depth $k$ is
$K_k^\mu$, and $E_k$ is that kernel dressed by $\nu$ (the Gram
$B B^T$ of~\cite{r283,r285}).  The identity
\[
\lambda_{\max}(E_k)<1
\qquad\Longleftrightarrow\qquad
H_k(w)\succ 0
\qquad\Longleftrightarrow\qquad
h_0(w),\ldots,h_{k-1}(w)>0
\]
is the finite dictionary of lemma $L^*$ at depth $k$.

Round 380 records two kernel theorems, used here as black boxes.
\texttt{rankOne\_inertia\_antitone}: $\indm(A+vv^T)\le\indm(A)$.
\texttt{adaptive\_band\_from\_entry}: if $A_{n+1}=A_n+v_n v_n^T$
and $\indm(A_{n_0})\le 1$, then $\indm(A_n)\le 1$ for every
$n\ge n_0$.  The Jacobi count
\texttt{indNeg\_hankel\_eq\_neg\_pivot\_count} identifies
$\indm H_{n_0+1}$ with the number of negative prefixes among
$h_0,\ldots,h_{n_0}$.  The entry lemma is therefore a statement
about those prefixes, at a source-defined $n_0$.

\section{What the recurrence sees}\label{sec:rec}

The scaled signed Stieltjes chain on $(X,w)$ produces the Jacobi
coefficients $\alpha_k,\beta_k^2$ with $\beta_k^2=h_k/h_{k-1}$.
A sign change of $h$ is exactly a negative $\beta^2$.

\begin{prop}[Recursion, census]\label{prop:alpha}
On the flagship window $z=16$ ($N=184$) the source chain has
$\lvert\alpha_k\rvert\le 0.57$ through degree $N/2$ and the first
$\lvert\alpha\rvert>2$ occurs at the first negative pivot
$k=N$.  The matched scramble (seed $1$) stays mild through
degree $20$ and explodes at the first negative pivot
($\alpha_{21}=+3.23$), with $\beta^2_{20}<0$ and
$\lambda_{\max}(E_{22})>1$.  The first dangerous approach
$h_k\to 0$ is not a source/scramble separator: both worlds drop
eight log-units of $\lvert h\rvert$ by degree $6$.
\end{prop}

The $\alpha$-explosion is a \emph{symptom} of the first sign
change, not a cause that could be bounded a priori from the
three-term coefficients alone.  Leg~A therefore kills the hope
that a sign/size profile of $(\alpha,\beta^2)$ without the
atomic support would distinguish the source: any bound that
uses only the recurrence until it breaks is a restatement of
the first negative pivot.

\section{Finite identities}\label{sec:id}

\begin{lemma}[RKHS sum]\label{lem:rkhs}
For every real polynomial $p$ of degree $<k$,
\[
\int p^2\,d\nu
\le
\Bigl(\sum_j \nu_j\,K_k^\mu(y_j,y_j)\Bigr)
\int p^2\,d\mu.
\]
Consequently $\sum_j\nu_j K_k^\mu(y_j,y_j)<1$ implies
$H_k(w)\succ 0$.
\end{lemma}

\begin{proof}
$p(y)=\langle p,K_k^\mu(\,\cdot\,,y)\rangle_\mu$ on the span of
degree $<k$, so $p(y)^2\le K_k^\mu(y,y)\int p^2\,d\mu$ by Cauchy--Schwarz
in $L^2(\mu)$.  Sum against $\nu$.
\end{proof}

The sum bound is SATZ and source-defined.  It is also coarse:
on the flagship window the largest $k$ with the sum strictly
less than one is $K_{\mathrm{sum}}=9=0.049\,N$, and the scramble
already has sum $1.78$ at $k=9$.  It does not yield
$K/N\to\kappa>0$.

\begin{lemma}[Pair energy]\label{lem:energy}
If $m_r=\int x^r\,dw$, then
\[
m_0 m_2-m_1^2
=\tfrac12\sum_{i,j}w_i w_j(x_i-x_j)^2.
\]
In particular $h_1>0$ whenever $m_0>0$ and the right-hand side is
positive.  On the five-atom interlacing toy of
Lemma~\ref{lem:toy-int} both sides equal $56$.
\end{lemma}

\begin{proof}
Expand the right-hand side:
$\sum_{i,j}w_i w_j(x_i^2-2x_i x_j+x_j^2)
=2m_0 m_2-2m_1^2$.
\end{proof}

\begin{lemma}[Christoffel comparison]\label{lem:christoffel}
Let $p$ run over monic polynomials of degree $k$.  Then
\[
h_k(w)
=\min_p\int p^2\,d(\mu-\nu)
\ge
\bigl(1-\lambda_{\max}(E_{k+1})\bigr)\,h_k(\mu),
\]
whenever $\lambda_{\max}(E_{k+1})$ is defined.  Since $h_k(\mu)>0$,
the sign of $h_k(w)$ is decided by whether $E_{k+1}$ is a strict
contraction --- the $L^*$ dictionary at depth $k+1$.  On the
three-atom flank toy of Lemma~\ref{lem:toy-flank} one has
$h_1(w)=6\ge(1-\tfrac13)\cdot 6=4$.
\end{lemma}

\begin{proof}
Degree $k$ is degree $<k+1$, so
$\int p^2\,d\nu\le\lambda_{\max}(E_{k+1})\int p^2\,d\mu$.
Take the $\mu$-minimizer of the monic class.
\end{proof}

\section{Interlacing toys}\label{sec:toy}

\begin{lemma}[Five-atom $2$-versus-$1$]\label{lem:toy-int}
On $\{-2,-1,0,1,2\}$ with weights $\{2,-1,2,-1,2\}$ one has
$h_0=4$, $h_1=14$, $h_2=13$, $h_3=-144/7$.  Here $S=5$, $N=3$,
$S_+=3$, $S_-=2$, and $h_0,h_1,h_2>0$ is the full free window of
the toy.  The first negative pivot is post-cap.
\end{lemma}

\begin{lemma}[Three-atom flank]\label{lem:toy-flank}
On $\{0,1,2\}$ with weights $\{3,-2,3\}$ the unique negative
run has flanking ratio $2/(3+3)=1/3$, and
$h_0=4$, $h_1=6$, $h_2=-3$.  Here $N=2$, and $h_0,h_1>0$ is
again the free window.  The one-atom kernel identity
$\nu(\{1\})K_2^\mu(1,1)=1/3=\lambda_{\max}(E_2)$ is exact.
\end{lemma}

\begin{lemma}[Equal two-period kills $H_1$]\label{lem:adv-tp}
On $\{0,1,2,3,4,5\}$ with weights $\{+1,-1\}^3$ one has
$H_1=0$.  Total mass vanishes, so $h_0=0$.
\end{lemma}

\begin{lemma}[Clustered run of three]\label{lem:toy-cluster}
On $\{0,\ldots,6\}$ with weights
$\{2,2,-1,-1,-1,2,2\}$ one has $H_1>0$, $H_2>0$, $H_3<0$.
Here $N=4$, and the first negative pivot sits at degree $2<N-1$.
A negative run of length $3$ is therefore not forbidden from
killing the free window by the moment count alone.
\end{lemma}

\section{The flank condition}\label{sec:flank}

Order the cosine nodes as a path.  A \emph{negative run} is a
maximal consecutive block of $\nu$-atoms.

\begin{definition}[Flank]\label{def:flank}
A signed atomic measure on a linearly ordered support is
\emph{$c$-flanked} if every negative run has length at most $2$
and the $\nu$-mass of the run is at most $c$ times the sum of
the at most two immediately adjacent $\mu$-atoms.
Write (F1) for the run-length bound, (F2) for $c$-flank with
$c<1$, and (F2$_{2/3}$) for $c=2/3$.
\end{definition}

The constant $2/3$ is the first rational strictly above every
core-$42$ flanking ratio (all $\le 0.619$).  It is not fitted to
the extension anchors; seven deep EXT windows have $2/3<r_{\max}<1$
and are excluded from (F2$_{2/3}$) while still satisfying (F1)
and (F2).  The $3/8$-floor and MED-CAP of the gap process
(r361) enter only as $\tau$-separation of the nodes: consecutive
cosine gaps are comparable, so a run of length $\le 2$ is a
$O(1/S)$ block in $\theta$.  They are not themselves a pivot
theorem.

\begin{prop}[Flank census]\label{prop:census}
On the sealed $85$-window MAIN ladder: (F1) holds $85/85$, (F2)
holds $85/85$ (every $r_{\max}<1$, no unflanked run), and
(F2$_{2/3}$) holds $78/85$.  The seven EXT-heavy windows are
$kz\in\{69,96,97,99,107,117,129\}$, with worst $r_{\max}=0.879$
at $kz=129$.  On the core $42$, no run has ratio $\ge 2/3$, and
at most $4.02\%$ of runs have ratio $\ge 1/2$.  The six
terminal-dead $\chi$-sprouts (chi3 $15,19,23,33,39$ and chi4
$20$) all satisfy (F1) and (F2$_{2/3}$); terminal death is not a
prefix-pivot failure.  The matched scramble fails (F1)
($\max$-run $5$, nine runs of length $\ge 3$) and (F2)
($r_{\max}=2.71$, twelve runs with ratio $\ge 1$), and its
first negative pivot is $21<\lfloor 2\cdot 184/5\rfloor=73$.
\end{prop}

Uniform four-node blocks with $\nu/\mu<1$ do \emph{not} hold on
MAIN (flagship overlapping $r_{\max}=2.08$).  The source
structure is the \emph{run} against its flanks, not a dyadic
partition of the grid.  Digamma/tent regularity of the weights
is used only through this discrete comparison: the same cosine
mesh with scrambled signs produces long runs and flank ratios
above one.

\section{Two-period amplitudes}\label{sec:tp}

\begin{prop}[Two-period table]\label{prop:twoperiod}
On the cosine grid of $S=81$ nodes ($N=41$), the signed two-period
$w_j=(1-x_j)(\pi/S)\cdot(+1$ even $j$, $-c$ odd $j)$ has first
negative pivot
\[
\begin{array}{c|cccc}
c & 1/2 & 2/3 & 4/5 & 1\\
\hline
\text{first neg.} & 40 & 20 & 10 & 0
\end{array}
\]
so $\mathrm{first}/N \approx 0.98,\;0.49,\;0.24,\;0$.  The same
ratios, to two digits, reproduce at $S=21,41,161$.  Amplitude
$c=2/3$ is therefore a half-filling adversary for a theorem that
uses only (F1) and (F2$_{2/3}$): it forbids claiming
$n_0=N-1$, and it permits $n_0=\lfloor 2N/5\rfloor=16$ on this
row ($20\ge 16$).  Amplitude $c\ge 1$ kills at $h_0$, matching
Lemma~\ref{lem:adv-tp}.
\end{prop}

A two-period of amplitude $0.88$ (the EXT-heavy $r_{\max}$) dies
at $\sim 0.12\,N$, which is the scramble wall.  That is why a
uniform-$c$ theorem covering the seven EXT-heavy windows cannot
beat the generic bound, and why (F2$_{2/3}$) is the honest
sufficient condition: it excludes those seven from the
$\kappa=2/5$ claim while still covering the core $42$ and the
dead $\chi$.  The seven remain under (F1)+(F2) with $c<1$, which
is enough for $h_0,h_1$ and for the RKHS-sum entry, not for a
positive fraction of $N$.

\section{The reduction}\label{sec:reduction}

\begin{definition}[Source-defined entry]\label{def:n0}
On a window of half-filling $N$ set
\[
n_0(w)=\lfloor \tfrac25 N\rfloor
\]
whenever $w$ is $(2/3)$-flanked in the sense of
Definition~\ref{def:flank}, and declare the entry failed
otherwise.  This $n_0$ is a function of the ordered signed
weight alone --- no spectral input, no first-negative readout.
\end{definition}

\begin{prop}[Entry on the census]\label{prop:entry}
On every MAIN window, $h_k>0$ for all $k\le n_0(w)$ (in fact
through $k=N-1$, the r377 free-window census).  On the scramble,
$n_0=73$ and the first negative pivot is $21<73$, so the
scramble is excluded by the premise of Definition~\ref{def:n0}
and by the conclusion.  On every two-period row of
Proposition~\ref{prop:twoperiod} at $c=2/3$, the first negative
pivot is at least $n_0$.
\end{prop}

The reduced lemma is therefore:

\begin{quote}
\emph{If $w$ is $(2/3)$-flanked, then $h_0,\ldots,h_{n_0}>0$ with
$n_0=\lfloor 2N/5\rfloor$.  Equivalently
$\indm H_{n_0+1}(w)=0\le 1$.  By
\texttt{adaptive\_band\_from\_entry} this entry activates the
rank-one band from $n_0$.}
\end{quote}

The implication from (F1)+(F2$_{2/3}$) to the prefix bound is
SATZ on the toys of~\S\ref{sec:toy}, on the two-period family
at $c=2/3$ (Proposition~\ref{prop:twoperiod}), and as a census
on the construction.  A cofinal $L^2$ proof of the same
implication is the named remainder of the next section.

\section{What half-filling still needs}\label{sec:half}

The RKHS sum (Lemma~\ref{lem:rkhs}) gives $K/N\to 0$.  An
$\lVert p\rVert_\infty$ plus Bernstein-in-$\theta$ estimate
grows like $\sqrt{S}$ and likewise cannot yield a uniform
positive fraction of $N$.  The bound that \emph{would} give
$K/N\to\kappa>0$ is an operator-norm estimate on flank blocks:
if $E_k$ were localized at scale $S/k$, then
$\lambda_{\max}(E_k)\le c+\mathrm{tail}_k$ with $c\le 2/3$, and
$\mathrm{tail}_k$ small for $k\le (2/5)N$.  At $k=N$ the same
estimate becomes $\lambda_{\max}(E_N)<1$, which is $L^*$.

Thus:

\begin{itemize}
\item $n_0=\lfloor 2N/5\rfloor$ under (F2$_{2/3}$) is the reduced
entry (strictly smaller than $L^*$, strictly larger than the
scramble death $0.114\,N$, strictly below the two-period kill
$\sim N/2$).
\item $n_0=N-1$ with zero sign changes is $L^*$ and is
\emph{not} a consequence of (F1)+(F2$_{2/3}$) alone: the
two-period of amplitude $2/3$ is a named adversary against
that claim.
\item The missing input for half-filling is the $L^2$
Christoffel tail on the flank blocks (r285 coherence, the
assist of $\lambda_{\max}=\mathrm{maxdiag}\cdot(1+\mathrm{assist})$),
together with whatever stronger source regularity (the
$\delta_{1/2}\le 0.04$ quantile of the core, the tent profile)
upgrades $c=2/3$ toward $c=1/2$, where the two-period table
reaches $n_0=N-2$.
\end{itemize}

\section{Kill battery}\label{sec:kill}

\begin{center}
\begin{tabular}{@{}llll@{}}
\toprule
world & (F1) & (F2$_{2/3}$) & first neg.\\
\midrule
MAIN $85$ & $85/85$ & $78/85$ & $\ge N$ (r377)\\
core $42$ & $42/42$ & $42/42$ & $\ge N$\\
dead $\chi$ (6) & $6/6$ & $6/6$ & $\ge n_0$\\
scramble w9 & fail (run $5$) & fail ($r_{\max}=2.71$) & $21<73$\\
two-period $c=2/3$ & pass & pass & $\sim N/2\ge n_0$\\
two-period $c=1$ & pass & fail & $0$\\
clustered run $6$ & fail & --- & $<n_0$\\
\bottomrule
\end{tabular}
\end{center}

Scramble drops at both premises of the sufficient condition
(named: run length and flank ratio).  Dead $\chi$ satisfy both;
(F2$_{2/3}$) is not the terminal-death mechanism.  The
two-period of amplitude $2/3$ is the binding adversary against
raising $\kappa$ from $2/5$ to $1$, and of amplitude $1$ against
dropping the $c<1$ clause.

\section{Lean docking (named Prop, not a sorry)}\label{sec:lean}

No Lean file is edited in this round.  The docking statement for
R384, against the r380 kernel in \texttt{RH/PivotCoordinate.lean},
is the following named Prop.  It is not a \texttt{sorry}: the
same convention as \texttt{ComplementaryDualHankelInertia} /
\texttt{CauchyInterlace}.

\begin{quote}
\texttt{def FlankEntryPrefix : Prop :=}\\
\quad\texttt{$\forall$ \{S : Nat\} (m : SignedAtoms S),}\\
\quad\texttt{\quad FlankTwoThird m $\to$}\\
\quad\texttt{\quad let N := (S+1)/2}\\
\quad\texttt{\quad let n0 := (2*N)/5}\\
\quad\texttt{\quad ($\forall$ i $\le$ n0, 0 $<$ m.h i)}
\end{quote}
Here $(2\cdot N)/5$ is Lean natural division, hence
$\lfloor 2N/5\rfloor$.  The previous real-number form
$2N/5\le n_0$ would have been false at non-multiples of $5$.

\texttt{FlankTwoThird} is Definition~\ref{def:flank} with
$c=2/3$ on the ordered support of \texttt{m}.  Given that Prop,
\texttt{indNeg\_hankel\_eq\_neg\_pivot\_count} turns the prefix
into $\indm H_{n_0+1}=0$, and \texttt{adaptive\_band\_from\_entry}
transports the bound along any Loewner chain issuing from that
entry.  The Prop does not assert (P1) or (P2) on a window, does
not assert $L^*$, and does not mention the Riemann hypothesis.

\section*{Claim boundary}

Finite identities, a named sufficient condition, a construction
census, and a reduced entry $n_0=\lfloor 2N/5\rfloor$.  Not a
theorem of $L^*$, not a theorem of the wall, and not evidence
for or against the Riemann hypothesis in either direction.

\begin{thebibliography}{9}
\bibitem{r283}
TFPT RH program, round 283: full-source quasi-definiteness and
the $L^*$ reduction.
\bibitem{r285}
TFPT RH program, round 285: Christoffel decomposition
$\lambda_{\max}=\mathrm{maxdiag}\cdot(1+\mathrm{assist})$.
\bibitem{r377}
S.~Hamann, \emph{Post-cap pivots of the signed source},
\texttt{rh/problem/postcap\_pivots.pdf}, 2026.
\bibitem{r380}
\texttt{RH/PivotCoordinate.lean}, round 380: rank-one inertia
and adaptive band, sorry-free.
\end{thebibliography}

\end{document}
